\section{Partie théorique}

\subsection{Oscillations libres d’un pendule de Pohl}

Le pendule de Pohl est un oscillateur mécanique rotatif composé d'un disque couplé à un ressort
spiral. La coordonnée généralisée du système est l’angle de torsion $\theta(t)$.
En appliquant le théorème du moment cinétique au disque, on obtient :

\[
\frac{dL}{dt} = \sum M_{\mathrm{ext}},
\qquad L = J\,\dot{\theta},
\]

où $J$ est le moment d’inertie du balancier. Deux moments mécaniques agissent sur le système :
(i) le moment de rappel du ressort spiral
\[
M_r = -C\,\theta,
\]
où $C$ est la constante de torsion, et (ii) le moment de frottement visqueux
\[
M_f = -b\,\dot{\theta}.
\]

En combinant ces contributions, l'équation du mouvement devient :

\[
J\ddot{\theta} + b\dot{\theta} + C\theta = 0.
\]

On définit la pulsation propre non amortie :
\[
\omega_0 = \sqrt{\frac{C}{J}},
\]
ainsi que le coefficient d’amortissement :
\[
\lambda = \frac{b}{2J}.
\]

L’équation différentielle se met sous la forme caractéristique :
\[
\ddot{\theta} + 2\lambda \dot{\theta} + \omega_0^2 \theta = 0.
\]

Dans le régime faiblement amorti ($\lambda < \omega_0$), la solution générale est :

\[
\theta(t) = \theta_0\,e^{-\lambda t}
\cos(\omega_1 t + \phi),
\]

où la pseudo-pulsation est :

\[
\omega_1 = \sqrt{\omega_0^2 - \lambda^2}.
\]

La période mesurée expérimentalement est alors :

\[
T = \frac{2\pi}{\omega_1}.
\]

L’expression de $J$ peut être inversée en fonction des paramètres mesurés :

\[
J = \frac{C}{\omega_0^2}.
\]


\subsection{Oscillations forcées}

Lorsque le moteur excite le système, un moment sinusoïdal est appliqué :

\[
M_e(t) = M_0 \cos(\omega t),
\]

ce qui modifie l'équation différentielle :

\[
J\ddot{\theta} + b\dot{\theta} + C\theta = M_0\cos(\omega t).
\]

En divisant par $J$, on obtient l'équation caractéristique sous forme normalisée :

\[
\ddot{\theta} + 2\lambda \dot{\theta} + \omega_0^2 \theta
= A\cos(\omega t),
\qquad A = \frac{M_0}{J}.
\]

\subsection{Méthode des phaseurs}

On cherche une solution sinusoïdale particulière de la forme :

\[
\theta(t) = \Theta(\omega)\cos(\omega t + \phi(\omega)).
\]

En utilisant la méthode complexe (phaseurs), l'amplitude en régime forcé est :

\[
\Theta(\omega)
= \frac{A}{\sqrt{(\omega_0^2 - \omega^2)^2 + (2\lambda\omega)^2}}.
\]

La phase du mouvement par rapport à la force excitatrice est :

\[
\phi(\omega)
= \arctan\!\left(\frac{2\lambda\omega}{\omega_0^2 - \omega^2}\right).
\]

La résonance apparaît lorsque l’amplitude est maximale, ce qui se produit pour :

\[
\omega_{\mathrm{res}} = \sqrt{\omega_0^2 - 2\lambda^2}.
\]

\subsection{Facteur de qualité}

Le facteur de qualité d’un oscillateur amorti est défini par :

\[
Q = \frac{\omega_0}{2\lambda}.
\]

Dans l’analyse fréquentielle (courbe de puissance), il peut aussi être obtenu via :

\[
Q = \frac{f_0}{\mathrm{FWHM}},
\]

où $\mathrm{FWHM}$ est la largeur à mi-hauteur du pic de résonance,
et $f_0 = \omega_0 / (2\pi)$ la fréquence propre.

